% =============================================================================
% Section 3.3: Background on AI in E-Learning
% =============================================================================

\section{Background on AI in E-Learning}
\label{sec:ai-elearning}

\subsection{OpenAI Whisper for Speech Recognition}

OpenAI Whisper~\cite{radford2023whisper} is a general-purpose speech recognition model trained on 680,000 hours of multilingual and multitask supervised data collected from the web. Unlike prior automatic speech recognition (ASR) systems that relied on meticulously curated, single-domain datasets, Whisper leverages weak supervision at an unprecedented scale, resulting in robust performance across diverse acoustic environments, speaker accents, and recording qualities. The model architecture is a sequence-to-sequence Transformer encoder-decoder: audio is resampled to 16 kHz, converted to an 80-channel log-mel spectrogram, and encoded by the Transformer encoder; the decoder then autoregressively produces the transcript text, optionally including timestamps, language identification, and segment boundaries.

Whisper's relevance to e-learning platforms is established by Panthel et al.~\cite{panthel2024whisper}, who evaluated Whisper transcriptions against human-produced subtitles across 120 university lecture recordings. They reported word error rates between 5.2\% and 8.1\% on English lectures with clear audio, rising to 12--15\% on recordings with background noise, accented speech, or specialised technical terminology. Critically, the study concluded that Whisper's output was ``sufficient for automated downstream tasks such as keyword extraction, content indexing, and AI quiz generation [...] with manual correction recommended only for public-facing subtitles where near-perfect accuracy is required.''

NexaLearn integrates Whisper through a dedicated \textbf{AI Pipeline} bounded context. When a video asset completes transcoding (Section~\ref{sec:video-streaming}), the pipeline:

\begin{enumerate}
  \item Downloads the MP4 rendition to a transient processing environment.
  \item Extracts the audio track and segments it into 30-second chunks aligned with the Whisper context window.
  \item Sends each chunk to the Whisper API or a self-hosted Whisper instance (configurable per deployment), requesting transcripts with word-level timestamps and segment boundaries.
  \item Assembles the per-chunk transcripts into a monolithic transcript document, preserving the temporal mapping between transcript segments and video offset.
  \item Stores the transcript as a JSON document associated with the \texttt{VideoAsset} entity, enabling downstream consumers to retrieve ``the transcript segment starting at second 145'' for fine-grained synchronisation.
\end{enumerate}

\subsection{LLM-Based Quiz Generation}

The second AI capability of NexaLearn is automated quiz generation from lecture transcripts using large language models. Kumar et al.~\cite{kumar2023llmquiz} demonstrated an end-to-end pipeline that inputs a lecture transcript and outputs a structured set of multiple-choice questions, true-or-false statements, and short-answer prompts, each with a correct answer, plausible distractors, and a justification. Their evaluation across 50 university lecture transcripts showed that GPT-4--generated questions were rated by domain experts as acceptable or better in 73\% of cases, with the primary failure mode being questions that tested recall of peripheral transcript details rather than conceptual understanding.

NexaLearn's AI pipeline generalises this approach. Upon receiving a transcript (either generated by Whisper or uploaded manually by the instructor), the pipeline constructs a prompt instructing the LLM to produce structured JSON adhering to a predefined schema:

\begin{lstlisting}[caption=AI-generated quiz JSON schema,label=lst:ai-quiz-schema]
{
  "questions": [
    {
      "stem": "What is the primary invariant enforced by ...?",
      "type": "multiple-choice",
      "options": [
        { "id": "A", "text": "...", "isCorrect": false },
        { "id": "B", "text": "...", "isCorrect": true  },
        { "id": "C", "text": "...", "isCorrect": false },
        { "id": "D", "text": "...", "isCorrect": false }
      ],
      "justification": "..."
    }
  ]
}
\end{lstlisting}

The generated quiz payload is not published directly to learners. Instead, NexaLearn applies a \textbf{human-in-the-loop review} workflow: the quiz is stored with a \texttt{READY\_FOR\_REVIEW} status, and the instructor is notified (via the Notifications context) to review, edit, approve, or reject each proposed question. Only upon instructor approval does the quiz status transition to \texttt{PUBLISHED}, making it available to enrolled learners. This review window ensures that the instructor retains pedagogical control while benefiting from automated question drafting---reducing the time required to create a quiz from a 60-minute lecture from approximately two hours of manual authoring to fifteen minutes of review and refinement.

\subsection{HMAC-Signed Callback Authentication}

The AI pipeline operates asynchronously: NexaLearn submits a transcription or quiz-generation request, the external AI service processes it (potentially taking seconds to minutes), and the service calls back a NexaLearn endpoint with the results. To authenticate these callbacks, NexaLearn implements an HMAC-signature scheme inspired by Stripe webhook signatures~\cite{stripe2024api}. Each callback request includes an \texttt{X-Pipeline-Signature} header with format \texttt{t=<unix-timestamp>,v1=<hex-digest>}. The recipient---the AI Pipeline controller---recomputes the HMAC-SHA256 digest of the request body using the shared secret, verifies that the timestamp is within a configurable tolerance window (default 300 seconds), and rejects requests whose signature does not match or whose timestamp is stale. This prevents replay attacks and ensures that only callbacks from the genuine AI pipeline are accepted.

The HMAC callback pattern, combined with the outbox event processing described in Section~\ref{sec:event-driven-outbox}, guarantees that AI-generated content is durably stored even if the application crashes during callback handling. The callback handler inserts the generated transcript or quiz into the database within a transaction that also creates an outbox event; the outbox processor subsequently notifies the instructor and updates related projections.
